\chapter{文档范围与实现基线}
\section{阅读目标与依据}
本文按现有工程梳理实现，面向维护、扩展、缺陷定位和交接。源码基线为提交\code{570e669}，应用版本0.40.0。文档日期采用Asia/Singapore时区。版本记录只用于解释背景；当记录与代码不一致时，以当前函数行为为准，并注明差异。

本册解释运行结构、职责与约束；配套《详细设计》给出字段、接口和计算。附录从源码提取定位，数值从TS和JSON生成；“建议”段才是后续方案。核心约束是失败不部分扣费、旧快照不追改、回放不发奖、车辆不重复使用。

\section{产品边界}
当前交付物是Windows x64原生单机策略游戏，包含基地、独立车间、科研、六格编队、普通战役、核心副本、世界资源点、指挥官和战报。世界据点为本地NPC，突袭不代表真人联机。金币与VIP为本机模拟，没有支付网关和在线账户服务。

工程有两个前端上下文。正式原生版由Godot和Node组成；历史网页原型使用React、Worker和IndexedDB。二者复用规则，但启动、通信和持久化适配不同。不能把浏览器事务、localStorage和原生HTTP服务写成同一个部署模型。
\src{README.md}{产品范围；从源码运行}

\section{核心概念与计量}
\begin{small}\begin{longtable}{@{}P{27mm}P{121mm}@{}}
\toprule 概念 & 工程含义\\\midrule\endhead
状态与视图 & \code{GameState}保存事实；\code{view()}派生报价、可用编队、调度和属性。视图不是第二个可写数据库。\\
命令与请求 & \code{Command}改变游戏状态；RPC的\code{op}还覆盖启动、读报告、导入导出及管理。二者并非一一对应。\\
工位与订单 & 工位是执行通道；\code{Job}可以包含多辆车并逐辆交付；\code{jobBacklog}保存等待订单。\\
编队与库存 & 编队保存计划，不预占车辆；远征出发才扣待命库存，并保存战斗与载重快照。\\
战斗与表现 & Node先结算整场\code{BattleReport}，Godot再播放事件，动画不驱动伤害结算。\\
时间与比例 & 时间采用毫秒整数，速率多用每小时；战斗比例通常用基点，10000等于100\%；经济函数可能返回百分数。\\
等级 & 建筑、玩家科研与经验等级最高120；统率、声望、战斗技能没有玩法上的120级封顶，仍受安全整数和存档规模约束。\\
\bottomrule\end{longtable}\end{small}
